\documentclass[../main.tex]{subfiles}

\begin{document}
    \subsection{\Pointer}
    {\Pointer} dapat didefinisikan sebagai suatu variabel yang menyimpat alamat memori. Jika kita mempunyai variabel dengan tipe data tertentu, maka untuk mendapatkan alamat dari variabel tersebut adalah dengan menggunakan operator \inputscript{\&} \parencite{raharjo2015pemrograman}.

    Kita dapat mendefinisikan sebuah pointer ke tipe data dengan cara mendeklarasikannya dalam bentuk \inputscript{*d}, dimana \inputscript{d} merupakan nama dari variabel yang didefinisikan. Simbol \inputscript{*} harus diulangi untuk setiap variabel {\pointer} \parencite{Lippman_Lajoie_Moo_2013}:

\begin{lstlisting}
 int *ip1, *ip2;  // both ip1 and ip2 are pointers to int
 double dp, *dp2; // dp2 is a pointer to double; dp is a double
\end{lstlisting}

    Jika kita mempunyai sebuah variabel yang bertipe \inputscript{long} (misalnya \inputscript{X}), maka kita dapat memerintahkan {\pointer} \inputscript{P} untuk menunjuk ke alamat yang ditempati oleh variabel \inputscript{X}. Untuk melakukan hal tersebut kita perlu menuliskan kode seperti berikut \parencite{raharjo2015pemrograman}:

\begin{lstlisting}
 // Mendeklarasikan variabel X dengan tipe long
 long X;
 // Mendeklarasikan pointer P
 long *P;

 // Memerintahkan P untuk menunjuk alamat
 // dari variabel X
 P = &X;
\end{lstlisting}

    \subsection{Referensi}

    Sebuah \textit{reference} merupakan sebuah alias atau sinonim untuk variabel lain. \textit{Reference} dideklarasikan dengan \textit{syntax} \inputscript{type\& ref\_name = var\_name;} dimana \inputscript{type} merupakan tipe data variabel, \inputscript{ref\_name} merupakan nama dari \textit{reference}, dan \inputscript{var\_name} merupakan nama dari variabel. Sebagai contoh \parencite{hubbard2021programming} :

\begin{lstlisting}
 int& rn = n; // r is a synonym for n
\end{lstlisting}

    Variabel \inputscript{rn} dideklarasikan sebagai sebuah \textit{reference} ke variabel \inputscript{n}, yang harus telah terdeklarasi.

\end{document}
